% In this paper audio signals are modelled with a generative model operating directly on the raw audio waveform.  
In this paper we introduce a new generative model operating directly on the raw audio waveform.  
The joint probability of a waveform $\vec{x} = \{ x_1, \dots, x_T \}$ is factorised as a product of conditional probabilities as follows:
\begin{equation}
p\left(\vec{x}\right) = \prod_{t=1}^{T} p\left(x_t \mid x_1, \dots ,x_{t-1}\right)
\label{eq:px}
\end{equation}
Each audio sample $x_t$ is therefore conditioned on the samples at all previous timesteps.

Similarly to PixelCNNs \citep{van2016pixel, ConditionalPixelCNN}, the conditional probability distribution is modelled by a stack of convolutional layers. There are no pooling layers in the network, and the output of the model has the same time dimensionality as the input. The model outputs a categorical distribution over the next value $x_{t}$ with a softmax layer and it is optimized to maximize the log-likelihood of the data w.r.t. the parameters. Because log-likelihoods are tractable, we tune hyper-parameters on a validation set and can easily measure if the model is overfitting or underfitting.

\subsection{Dilated Causal Convolutions}

\begin{figure}[h]
\centering
\includegraphics[width=0.85\linewidth]{Figures/masked_convolutions2.pdf}
\caption{Visualization of a stack of causal convolutional layers.}
\label{fig:masked_convolution}
\end{figure}

The main ingredient of WaveNet are causal convolutions. By using causal convolutions, we make sure the model cannot violate the ordering in which we model the data: the prediction $p\left(x_{t+1} \mid x_1,...,x_{t}\right)$ emitted by the model at timestep $t$ cannot depend on any of the future timesteps $x_{t+1}, x_{t+2},\dots,x_T$ as shown in \figref{fig:masked_convolution}. For images, the equivalent of a causal convolution is a masked convolution \citep{van2016pixel} which can be implemented by constructing a mask tensor and doing an elementwise multiplication of this mask with the convolution kernel before applying it. For 1-D data such as audio one can more easily implement this by shifting the output of a normal convolution by a few timesteps. 

At training time, the conditional predictions for all timesteps can be made in parallel because all timesteps of ground truth $\vec{x}$ are known. When generating with the model, the predictions are sequential: after each sample is predicted, it is fed back into the network to predict the next sample.

Because models with causal convolutions do not have recurrent connections, they are typically faster to train than RNNs, especially when applied to very long sequences. One of the problems of causal convolutions is that they require many layers, or large filters to increase the receptive field. For example, in \figref{fig:masked_convolution} the receptive field is only 5 (= \#layers + filter length - 1). In this paper we use dilated convolutions to increase the receptive field by orders of magnitude, without greatly increasing computational cost. 

A dilated convolution (also called \emph{\`a trous}, or convolution with holes) is a convolution where the filter is applied over an area larger than its length by skipping input values with a certain step. It is equivalent to a convolution with a larger filter derived from the original filter by dilating it with zeros, but is significantly more efficient. A dilated convolution effectively allows the network to operate on a coarser scale than with a normal convolution. This is similar to pooling or strided convolutions, but here the output has the same size as the input. As a special case, dilated convolution with dilation $1$ yields the standard convolution.  \figref{fig:masked_dilated_convolution} depicts dilated causal convolutions for dilations $1$, $2$, $4$, and $8$. Dilated convolutions have previously been used in various contexts, \eg signal processing \citep{Holschneider1989,Dutilleux1989}, and image segmentation \citep{chen14semantic,YuKoltun2016}.

\begin{figure}[ht]
\centering
\includegraphics[trim={6.625in 0 4.95in 0},clip,width=0.85\linewidth]{Figures/dilated_masked_convolutions4.pdf}
\caption{Visualization of a stack of \emph{dilated} causal convolutional layers.}
\label{fig:masked_dilated_convolution}
\end{figure}

Stacked dilated convolutions enable networks to have very large receptive fields with just a few layers, while preserving the input resolution throughout the network as well as computational efficiency. In this paper, the dilation is doubled for every layer up to a limit and then repeated: \eg $$1,2,4,\dots,512,1,2,4,\dots,512,1,2,4,\dots,512.$$
The intuition behind this configuration is two-fold. First, exponentially increasing the dilation factor results in exponential receptive field growth with depth~\citep{YuKoltun2016}.  For example each $1,2,4,\dots,512$ block has receptive field of size $1024$, and can be seen as a more efficient and discriminative (non-linear) counterpart of a $1\times1024$ convolution. Second, stacking these blocks further increases the model capacity and the receptive field size.

\subsection{Softmax distributions}

One approach to modeling the conditional distributions $p\left(x_t \mid x_1, \dots ,x_{t-1}\right)$ over the individual audio samples would be to use a mixture model such as a mixture density network \citep{MDN} or mixture of conditional Gaussian scale mixtures (MCGSM) \citep{theis2015generative}.
However, \cite{van2016pixel} showed that a softmax distribution tends to work better, even when the data is implicitly continuous (as is the case for image pixel intensities or audio sample values). One of the reasons is that a categorical distribution is more flexible and can more easily model arbitrary distributions because it makes no assumptions about their shape.

Because raw audio is typically stored as a sequence of 16-bit integer values (one per timestep), a softmax layer would need to output 65,536 probabilities per timestep to model all possible values. To make this more tractable, we first apply a $\mu$-law companding transformation \citep{G711} to the data, and then quantize it to 256 possible values:
$$
f\left(x_t\right) = \operatorname{sign}(x_t) \frac{\ln \left(1+\mu \left| x_t \right|\right)}{\ln \left(1+\mu\right)},
$$
where $-1 < x_t < 1$ and $\mu = 255$. 
This non-linear quantization produces a significantly better reconstruction than a simple linear quantization scheme. Especially for speech, we found that the reconstructed signal after quantization sounded very similar to the original.

\subsection{Gated activation units}

We use the same gated activation unit as used in the gated PixelCNN \citep{ConditionalPixelCNN}:
\begin{equation}
\vec{z} = \tanh \left(W_{f, k} \ast \vec{x}\right) \odot \sigma \left(W_{g, k} \ast \vec{x} \right), \label{eq:gated_activation}
\end{equation}
where $\ast$ denotes a convolution operator, $\odot$ denotes an element-wise multiplication operator, $\sigma(\cdot)$ is a sigmoid function, $k$ is the layer index, $f$ and $g$ denote filter and gate, respectively, and $W$ is a learnable convolution filter.
In our initial experiments, we observed that this non-linearity worked significantly better than the rectified linear activation function \citep{nair2010rectified} for modeling audio signals.

\subsection{Residual and skip connections}

\begin{figure}[h]
\centering
\includegraphics[width=0.75\linewidth]{Figures/architecture2.pdf}
\caption{Overview of the residual block and the entire architecture.}
\label{fig:architecture}
\end{figure}

Both residual \citep{he15deep} and parameterised skip connections are used throughout the network, to speed up convergence and enable training of much deeper models. In \figref{fig:architecture} we show a residual block of our model, which is stacked many times in the network.

\subsection{Conditional WaveNets}

Given an additional input $\vec{h}$, WaveNets can model the conditional distribution $p\left(\vec{x} \mid \vec{h}\right)$ of the audio given this input. \eqnref{eq:px} now becomes
\begin{equation}
p\left( \vec{x} \mid \vec{h} \right) = \prod_{t=1}^{T} p\left(x_t \mid x_1, \dots ,x_{t-1}, \vec{h}\right).
\label{eq:pxh}
\end{equation}

By conditioning the model on other input variables, we can guide WaveNet's generation to produce audio with the required characteristics. For example, in a multi-speaker setting we can choose the speaker by feeding the speaker identity to the model as an extra input. Similarly, for TTS we need to feed information about the text as an extra input.

We condition the model on other inputs in two different ways: global conditioning and local conditioning. Global conditioning is characterised by a single latent representation $\vec{h}$ that influences the output distribution across all timesteps, \eg a speaker embedding in a TTS model. The activation function from \eqnref{eq:gated_activation} now becomes:
$$
\vec{z} = \tanh \left(W_{f, k} \ast \vec{x} + V_{f, k}^T \vec{h}\right) \odot \sigma \left(W_{g, k} \ast \vec{x} + V_{g, k}^T \vec{h} \right). \label{eq:global_conditioning}
$$
where $V_{*, k}$ is a learnable linear projection, and the vector $V_{*, k}^T\vec{h}$ is broadcast over the time dimension.

For local conditioning we have a second timeseries $h_t$, possibly with a lower sampling frequency than the audio signal, \eg linguistic features in a TTS model. We first transform this time series using a transposed convolutional network (learned upsampling) that maps it to a new time series $\vec{y} = f(\vec{h})$ with the same resolution as the audio signal, which is then used in the activation unit as follows:
$$
\vec{z} = \tanh \left(W_{f, k} \ast \vec{x} + V_{f, k} \ast \vec{y}\right) \odot \sigma \left(W_{g, k} \ast \vec{x} + V_{g, k} \ast \vec{y}\right), \label{eq:local_conditioning}
$$
where $V_{f, k} \ast \vec{y}$ is now a $1\times 1$ convolution. As an alternative to the transposed convolutional network, it is also possible to use $V_{f, k} \ast \vec{h}$ and repeat these values across time. We saw that this worked slightly worse in our experiments.

\subsection{Context stacks}

We have already mentioned several different ways to increase the receptive field size of a WaveNet: increasing the number of dilation stages, using more layers, larger filters, greater dilation factors, or a combination thereof. A complementary approach is to use a separate, smaller \emph{context} stack that processes a long part of the audio signal and locally conditions a larger WaveNet that processes only a smaller part of the audio signal (cropped at the end). One can use multiple context stacks with varying lengths and numbers of hidden units. Stacks with larger receptive fields have fewer units per layer. Context stacks can also have pooling layers to run at a lower frequency. This keeps the computational requirements at a reasonable level and is consistent with the intuition that less capacity is required to model temporal correlations at longer timescales.
